\section{Background}
Decision tree is a simple and representative model of machine learning. Decision trees operate according to the principle of "divide-and-conquer"; through a series of questions, hereafter referred to as features, they attempt to predict certain attributes, or labels, of the input data. Decision tree are often built greedily by prioritizing questions that obtain the greatest difference among the partitions of its input. Numerous papers have shown ensemble decision tree models, meaning those that leverage multiple parallel decision trees, namely Random Forest (RF) and Gradient Boosting (GBoost) (or its variant, Extreme Gradient Boosting (XGBoost)), to outperform the rest.


Random forests, or RF, is an ensemble learning method that constructs multiple decision trees during training and outputs the mode of the classes (classification) or mean prediction (regression) of the individual trees. The algorithm synchronously generates multiple trees of a desired depth; randomness is introduced via feature selection for each tree's branches, as well as through the use of bootstrapped samples of the training data. This approach helps to reduce overfitting and improve generalization. RF is widely used for its robustness and ability to handle large datasets with high dimensionality.

Gradient boosting, or GBoost, is another ensemble learning method that builds decision trees sequentially, where each tree attempts to correct the errors of the previous one. The algorithm optimizes a loss function by adding new trees that minimize the residuals of the predictions made by the existing ensemble. Extreme Gradient Boosting (XGBoost) is a popular variant of GBoost that incorporates regularization techniques to prevent overfitting and improve performance. Both GBoost and XGBoost are known for their high accuracy and efficiency in handling complex datasets.

Another important aspect of decision tree-based models is the evaluation metrics used to assess their performance. Common metrics include accuracy, precision, recall, F1-score, and area under the receiver operating characteristic curve (AUC-ROC). These metrics provide insights into the model's ability to correctly classify instances, handle imbalanced datasets, and make informed decisions based on the predictions.

Finally, techniques for handling imbalanced datasets aer crucial when working with DT-based models to avoid bias towards the majority class in the models. Here, we focus on Synthetic Minority Over-sampling Technique (SMOTE), specifically its variant ADASYN, which utilies a weighted distribution to determine the learning difficulty of minority class samples and generates synthetic data accordingly. This approach helps to balance the dataset and improve the performance of the decision tree models in predicting the minority class.